% Source: https://artofproblemsolving.com/community/c3h1336325_quadratic_equation
% Stable identity: aops-topic-1336325-q01
Write $A=1-q+\dfrac{p^2}{2}$, $B=p(1+q)$, $C=q(q-1)+\dfrac{p^2}{2}$. Since the equation is quadratic, $A\ne0$, and it has a repeated root exactly when $B^2-4AC=0$.

Expand:
\[
4AC=\left(2-2q+p^2\right)\left(2q^2-2q+p^2\right)
=p^4+p^2\left(2q^2-4q+2\right)+\left(4q^2-4q^3-4q+4q^2\right)
\]
\[
=p^4+2p^2(q-1)^2-4q(q-1)^2 ,
\]
and $B^2=p^2(1+q)^2$. Hence
\[
B^2-4AC=p^2\left[(q+1)^2-2(q-1)^2\right]-p^4+4q(q-1)^2 .
\]
Since $(q+1)^2=(q-1)^2+4q$, the bracket is $4q-(q-1)^2$, and
\[
B^2-4AC=4p^2q-p^2(q-1)^2-p^4+4q(q-1)^2=-\left(p^2-4q\right)\left(p^2+(q-1)^2\right).
\]
(Check by expanding the right side: $-\left(p^4+p^2(q-1)^2-4p^2q-4q(q-1)^2\right)$.)

So $B^2-4AC=0$ means $p^2=4q$ or $p^2+(q-1)^2=0$. The second case forces $p=0$ and $q=1$, but then $A=1-1+0=0$ and the equation is not quadratic (all its coefficients vanish). Therefore $\boxed{p^2=4q}$.
